% Who, Not Why -- ICBINB-BIO @ NeurIPS 2026
% Compile on Overleaf: upload this file + neurips_2026.sty + the figures/ folder.
% Submission is DOUBLE-BLIND: keep the dblblindworkshop option, no author names,
% and point the repository link at an anonymised mirror before submitting.
\documentclass{article}
\usepackage[dblblindworkshop]{neurips_2026}
\workshoptitle{I Can't Believe It's Not Better (ICBINB): Failure Modes of AI in Biology}

\usepackage[utf8]{inputenc}
\usepackage[T1]{fontenc}
\usepackage{hyperref}
\usepackage{url}
\usepackage{booktabs}
\usepackage{amsfonts}
\usepackage{nicefrac}
\usepackage{microtype}
\usepackage{xcolor}
\usepackage{graphicx}

\title{Who, Not Why: Individual and Site Identity Dominate\\Frozen-Encoder Evaluation of Animal Vocalisations}

\author{%
  Anonymous Author(s)\\
  Affiliation\\
  \texttt{email}
}

\begin{document}
\maketitle

\begin{abstract}
Frozen audio encoders have become the default feature extractor for animal-vocalisation tasks, and
they are routinely evaluated with random train/test splits. We show this measures the wrong thing.
Across four public corpora---cat meows (21 individuals), dog barks (10), pig calls (6 recording labs)
and wild Egyptian fruit bats (10 emitters)---and six feature sets, individual and site identity are
far more decodable than behavioural context (identity 0.63--0.94 against chance 0.05--0.17), and
random splits inflate reported context accuracy by 10 to 24 points relative to held-out-group splits.
Within a corpus, the feature set that leaks less identity inflates less (partial $r{=}0.646$,
$p{=}7\times10^{-11}$); across corpora we do not detect this relationship ($r{=}0.221$, $n{=}12$).
Two consequences are concrete: on pig valence under held-out-lab evaluation, 88 classic acoustic
descriptors match a self-supervised speech encoder (0.669 vs.\ 0.660) while leaking less identity,
erasing a 6-point random-split lead; and the corpus's own published feature set collapses to chance
once a lab is held out (0.386 as distributed, 0.537 once the 700 clips its feature table does not
cover are excluded) against 0.575--0.745 on a random split. The failure extends to uncertainty quantification: conformal
prediction at $\alpha{=}0.10$ returns on-target marginal coverage (0.89--0.90) while the worst
individual receives 0.412, and the group-conditional remedy is unavailable by construction under a
group-disjoint split. We quantify the labelling budget that repairs each failure and recommend
reporting identity decodability alongside every context result. All experiments run on a CPU laptop.
\end{abstract}

\section{Introduction}

Two conventions have become standard in computational bioacoustics. First, use a large pretrained
audio encoder---BirdNET, Perch, AVES, or a speech model such as WavLM---frozen, as a feature
extractor. Second, evaluate the resulting classifier on a random split of the available clips.

The first convention is well justified. This paper is about the second. Animal-vocalisation corpora
are small and collected from a handful of individuals, often in a handful of sessions or labs. A
random split places clips from the same animal, recorded minutes apart in the same room, on both
sides of the train/test boundary. If the encoder represents \emph{who is calling} and \emph{where the
recording was made}---and we show it represents both very strongly---then a random-split evaluation
rewards recognising the animal and reports the result as understanding of its behaviour.

Readers from clinical and genomic ML will recognise the shape of this problem: a model that reads the
acquisition site rather than the biology, scored by a metric that cannot tell the difference
\citep{degrave2021,zare2026}.

That identity confounding inflates evaluation is not a new warning. It was named for biomedical data
by \citet{chaibubneto2019}, measured for audio by \citet{han2022}, and recommended against in
bioacoustics by \citet{colonna2016} and \citet{ghani2026}. What is missing is a measurement of
\emph{what the warning is worth} in this field. That is what we supply.

\paragraph{Contributions.}
(1)~A four-corpus, six-feature-set audit under one honest protocol (\S\ref{sec:audit}).
(2)~The identity--inflation relationship, scoped: within-corpus yes, cross-corpus not detected.
(3)~An unexpectedly strong simple baseline, and a published feature set that collapses to chance
under site shift.
(4)~A control showing a widely used dataset's context label is partly predictable from the
\emph{background} alone.
(5)~A channel stress test: the context axis is 1.4--4.2$\times$ more robust than identity
(\S\ref{sec:stress}).
(6)~The failure propagates to uncertainty quantification (\S\ref{sec:conformal}).
(7)~Recovery curves giving the labelling budget that repairs each failure (\S\ref{sec:recovery}).

\section{Related work}

\paragraph{Identity confounding outside bioacoustics.}
\citet{saeb2017} and \citet{chaibubneto2019} formalised record-wise versus subject-wise
cross-validation for digital health. \citet{han2022} measure it for audio health screening:
participant-independent AUC 0.71 against a random-split sensitivity/specificity of 0.84/0.78, with the
gain traced to overlapping participants. Speaker-independent evaluation has been standard in speech
paralinguistics since the INTERSPEECH challenge series \citep{schuller2009,schuller2010}, and
\citet{wagner2023} show the gap persists for modern self-supervised models. \citet{lin2026} report an
``identity trap'' in EEG foundation models, where variance attributable to subject identity dominates
task-related variance by one to two orders of magnitude; \citet{zare2026} document the analogous
dataset/site effect.

\paragraph{Shortcut learning.}
The general framing is \citet{geirhos2020}; the canonical scientific demonstration is
\citet{degrave2021}, where COVID-19 radiograph classifiers read laterality markers and hospital
source rather than pathology. \citet{kapoor2023} document leakage as a reproducibility crisis across
scientific ML. Our result is the bioacoustic instance.

\paragraph{Bioacoustic evaluation practice.}
\citet{colonna2016} recommended individual-disjoint splits for anuran calls; \citet{ghani2026} repeat
this as field guidance. \citet{stowell2019} used adversarial background mixing to expose
environmental confounding in bird classifiers. Benchmarks in current use \citep{hagiwara2023beans}
and the encoders evaluated on them \citep{hagiwara2023aves,miron2026avex} largely retain random or
pre-specified splits, and \citet{schwinger2026} show calibration varies sharply across datasets.
Notably, the ingredients of our claim already appear unremarked: \citet{abzaliev2024} report dog
individual-ID at 0.50 against 0.05 chance alongside context classification at 0.62 against a 0.56
baseline, without commenting on the disparity.

\paragraph{Conformal prediction.}
Split conformal prediction is due to \citet{vovk2013}; see \citet{angelopoulos2023} for a modern
treatment. \citet{barber2021} prove exact conditional coverage is unattainable in general; \citet{ding2023} give
the class-conditional sample-size bound $(1/\alpha)-1$ that we instantiate empirically in
\S\ref{sec:conformal}. To our knowledge there is no prior application of conformal prediction to bioacoustics or passive
acoustic monitoring.

\section{Protocol}

For every (dataset, feature set, layer) we report three numbers: \textbf{honest}---balanced accuracy
with an entire group held out; \textbf{leaky}---balanced accuracy under a random stratified 5-fold
split; and \textbf{identity}---accuracy of predicting the group label from the same features. The
third is the diagnostic that explains the gap between the first two.

Groups are the individual animal (cats, dogs, bats) or the recording lab (pigs). Site effects have a
concrete physical basis: \citet{turgeon2017} measured substantial microphone-to-microphone variation
between nominally identical autonomous recording units, and \citet{lostanlen2019} document the
resulting channel mismatch in deployed monitoring. Probes are
$\ell_2$-regularised logistic regression on standardised, mean-pooled embeddings, one layer at a
time; a linear probe measures what is linearly available, which is the setting practitioners use.
Feature sets: WavLM-base-plus, HuBERT-base, wav2vec2-base, AVES-bio, eGeMAPSv02 (88 openSMILE
functionals), and an 85-dimensional MFCC/log-mel summary. All frozen.

\section{The audit}\label{sec:audit}

\begin{table}[t]
\caption{Held-out-group vs.\ random-split evaluation, best layer per feature set. Chance is 0.333
(cats, dogs), 0.500 (pigs). Bats are reported for identity only: no bat context label is decodable.}
\label{tab:audit}
\centering\small
\begin{tabular}{llcccc}
\toprule
Dataset & Features & Held-out & Random & Inflation & Group identity \\
\midrule
Cats & WavLM & 0.559 & 0.725 & $+$0.167 & 0.793 \\
Cats & HuBERT & 0.571 & 0.671 & $+$0.099 & 0.700 \\
Cats & wav2vec2 & 0.554 & 0.737 & $+$0.183 & 0.755 \\
Cats & AVES-bio & 0.540 & 0.736 & $+$0.196 & 0.793 \\
Cats & eGeMAPS & 0.492 & 0.628 & $+$0.136 & 0.727 \\
Cats & MFCC & 0.399 & 0.602 & $+$0.203 & 0.768 \\
\midrule
Dogs & WavLM & 0.728 & 0.878 & $+$0.150 & 0.817 \\
Dogs & eGeMAPS & 0.619 & 0.733 & $+$0.114 & 0.792 \\
Dogs & MFCC & 0.551 & 0.793 & $+$0.242 & 0.892 \\
Dogs & duration only & 0.341 & 0.360 & $+$0.018 & 0.222 \\
\midrule
Pigs & WavLM & 0.660 & 0.883 & $+$0.223 & 0.938 \\
Pigs & eGeMAPS & \textbf{0.669} & 0.823 & $+$0.153 & 0.865 \\
Pigs & published 18 feats & 0.386 & 0.575 & $+$0.189 & 0.513 \\
Pigs & published 18 feats$^{\dagger}$ & 0.537 & 0.745 & $+$0.208 & 0.650 \\
\bottomrule
\end{tabular}

\vspace{2pt}
{\footnotesize $^{\dagger}$Same features and protocol, with the 700 clips omitted whose every
feature is absent from the published table (see \S\ref{sec:baseline}). Removing them raises the
held-out-lab score above chance and \emph{increases} the inflation.}
\end{table}

Inflation ranges from $+$0.099 to $+$0.242 and is present in every dataset and every feature set,
including the two hand-crafted ones. This is an evaluation pathology, not a neural-network pathology.

\subsection{Identity leakage predicts inflation---within a corpus}

Pooling all points, identity decodability correlates with the inflation gap at $r{=}0.611$
($n{=}83$, $p{=}9\times10^{-10}$); the within-dataset partial correlation is $r{=}0.646$
($p{=}7\times10^{-11}$). It survives restricting to layer-only variation inside a single encoder
($r{=}0.636$), holds in all six encoder clusters (sign test $p{=}0.031$), and is not a dimensionality
artifact (neural-only $r{=}0.694$). Identity tracks the \emph{leaky} number ($r{=}0.754$) more than
the honest one ($r{=}0.509$), as the leakage account requires.

\paragraph{The limit.} Across the 12 independent dataset\,$\times$\,feature-set units the correlation
is $r{=}0.221$ ($p{=}0.489$). At $n{=}12$ this test has roughly 15\% power, so we state that we
\emph{do not detect} a cross-corpus relationship, not that none exists.

\subsection{An unexpectedly strong simple baseline}\label{sec:baseline}

On pig valence under held-out-lab evaluation eGeMAPS scores 0.669 and WavLM 0.660, while eGeMAPS
leaks less identity (0.865 vs.\ 0.938). Under a random split WavLM appears clearly ahead (0.883 vs.\
0.823). We state this carefully: 0.9 points across four labs is not a significant win, and we do not
claim one. The defensible claim is that \textbf{the encoder's apparent advantage disappears under
honest evaluation while its identity leakage does not}. The ordering runs the other way in adjacent
work---\citet{tang2023} find SSL representations ahead of eGeMAPS in cross-language emotion
recognition---so this is domain-specific, not a general claim about feature families.

Separately, the 18 features published with the pig corpus score 0.386 under held-out-lab
evaluation---below the 0.500 chance level---against 0.575 on a random split. That below-chance figure
is not robust, and the reason is worth stating. The published feature table covers 4{,}331 of the
5{,}031 clips; for the remaining 700, every one of the 18 measurements is absent. All 700 are from a
single lab (NMBU, 51\% of its clips) and all carry the negative label, so under mean imputation they
become a single constant vector with one label attached. Excluding them, the same features and the
same protocol give 0.537 held-out-lab against 0.745 on a random split: at chance rather than below
it, and with inflation rising from $+$0.189 to $+$0.208. We report both rows. The claim the data
supports is that this feature set retains no cross-lab valence signal, not that its mapping inverts;
the leakage result, which is this paper's subject, is unaffected and slightly strengthened. Our
re-analysis uses a linear probe; the original study used other classifiers and did not claim cross-lab
generalisation. This is a reframing under a different protocol, not a refutation.

\subsection{The label can be predicted from the background}

Using non-vocal parts of a recording to expose confounding has precedent \citep{stowell2019}. We
apply it differently: probing behavioural context directly from background frames. Discarding the
vocalisation and keeping only the quietest 30\% of frames still yields 0.725 on the binary cat task.
The corpus documentation contains both halves of the problem: the protocol transfers cats to ``a room
in a different apartment or an office'' for the isolation condition, while the methods section states
room characteristics ``should not affect the captured audio signals'' \citep{ntalampiras2019}.

\section{Channel robustness}\label{sec:stress}

Perturbing the same clips and counting decision flips against probes trained on clean audio: at layer
9, context flips 1.1\% under a 4\,kHz low-pass against 4.7\% for a \emph{difficulty-matched} identity
probe (4.2$\times$); reverb 1.9$\times$; noise 1.4$\times$. Matching difficulty matters---the
unmatched comparison exaggerates the ratio to 2.7--6.1$\times$. Individual identity lives in fine
spectral detail, affect in coarse envelope structure. This also constrains the previous section: if
the context signal were purely room acoustics, reverb would scramble it, and it does not.

\section{The guarantees break too}\label{sec:conformal}

\begin{table}[t]
\caption{Split conformal at $\alpha{=}0.10$ (nominal coverage 0.90), calibrated and tested by group.}
\label{tab:conformal}
\centering\small
\begin{tabular}{lcccc}
\toprule
Dataset & Marginal & Worst group & IQR & Groups $<0.80$ \\
\midrule
Cats (20 individuals) & 0.904 & \textbf{0.412} & 0.186 & 5/20 \\
Dogs (10 individuals) & 0.892 & 0.624 & 0.059 & 2/10 \\
Pigs (6 labs) & 0.891 & 0.735 & 0.079 & 1/6 \\
\bottomrule
\end{tabular}
\end{table}

This is the conditional-coverage gap \citep{barber2021}, not an implementation error: a synthetic
exchangeable control returns 0.9048 over 2{,}000 trials and the per-group spread is equally wide under
an exchangeable split.

\paragraph{The remedy is unavailable.} Mondrian conformal conditioned on the group is bit-for-bit
identical to pooled under a group-disjoint split, because no test group has calibration data
(verified: 0/20, 0/10, 0/6 groups receive their own threshold). Conditioning on predicted class does
not help (worst cat $0.412\to0.425$). Given some of the deployment group's own labels, repair is
complete where data allows (dogs $0.696\to0.905$; pigs $0.775\to0.898$) but requires at least
$(1/\alpha)-1 = 9$ labelled clips from that group, which only 6 of 20 cats possess. \textbf{This
threshold is a theorem, not our finding}---\citet{ding2023} show the class-conditional quantile is
infinite below it, which makes Mondrian vacuity a corollary. Our contribution is measuring how often
bioacoustic corpora fall below the bound (14 of 20 cats) while marginal coverage still looks correct.

\paragraph{Shift.} Calibrating on dogs and testing on cats collapses even marginal coverage to 0.747.
Swapping individuals within a species costs nothing marginally; swapping species costs 15--35 points.

\section{How much labelling repairs it}\label{sec:recovery}

Adding $k$ labelled clips from the held-out group back into training: cats reach 0.727 at
$\approx$11 effective clips against a 0.729 random ceiling---the individual gap closes completely and
saturates, because the median cat contributes only $\approx$19 clips. Pigs do not saturate: 100 clips
from the target lab close 50.2\% of the gap and the curve is still climbing. On pigs eGeMAPS
\emph{starts} ahead (0.677 vs.\ 0.654 at $k{=}0$) but gains only $+$0.027 from 100 clips against
WavLM's $+$0.115: the encoder is more \emph{adaptable}, not more \emph{transferable}.

\section{Limitations}

One encoder family dominates; AVES is the only bioacoustic encoder tested at scale. Bats have no
decodable context label (0.301 vs.\ 0.250 chance) and are excluded from the inflation fit. Dog context
is confounded with recording session, so our dog numbers are leave-one-individual-out, not
leave-one-session-out. \textbf{The pig corpus confounds site with label distribution by design}: two
of six labs contribute a single valence class, the lab identifier alone predicts valence at 0.766 ---
above the 0.654 an encoder earns honestly --- and no behavioural context is shared across all six
labs, so leave-one-lab-out there conflates site shift with label-distribution shift and context
novelty. Only four labs are scorable and the baseline 95\% CI is [0.526, 0.725]. The pig numbers,
including the eGeMAPS comparison, should be read with that caveat. The cross-corpus correlation is underpowered ($n{=}12$). Probes are linear and
best layers are selected post hoc; we report selection-free means alongside. Extreme per-group
conformal minima come partly from small groups: restricted to cats with $n\geq10$ the worst is 0.748
and 3/15 fall below 0.80.

\section{Recommendations}

Split by group, never randomly. Report identity decodability next to every context number---it is one
extra probe. Report per-group spread, not only the mean, for accuracy and for coverage. Include a
hand-crafted baseline. State the target-group labelling budget if claiming a per-group guarantee.

\section*{Reproducibility and LLM use}

All code, cached results and figures are available at \textcolor{red}{[ANONYMISED MIRROR URL]}. Every
experiment runs on a CPU laptop; total compute is zero GPU-hours. Datasets are public: CatMeows \citep{ludovico2021,ntalampiras2019} (Zenodo 4008297, CC-BY-4.0), dog
barks \citep{molnar2008}, Soundwel \citep{briefer2022} (Zenodo 8252482, CC-BY-4.0), and Egyptian
fruit bats \citep{prat2017}. \textcolor{red}{[LLM-use disclosure required by the venue: state
here how LLM assistance was used in experimentation, analysis and writing.]}

\bibliographystyle{plainnat}
\bibliography{refs}

\end{document}
